
\section{Introduction}\label{sec:Introduction}

The availability of large amounts of multilingual corpus data has
opened up new possibilities for empirical approaches to contrastive
linguistics and translation and interpreting studies.
For example, we can empirically determine the contexts in which
comparable linguistic expressions from different languages are used
(e.g.
\citealt{Atayan2021a,Gast2015,Gast2017So,Gast2017Even,Gast2019,Gast2022}),
and we can investigate the conditions under which specific expressions
are translated in a specific way,
(e.g. \citealt{marco_2018,Cribleetal2019,DefrancqMagnifico2020,Defrancq2021,Atayan2021b,Yungetal2023}). These
possibilities come with new challenges, however. From the perspective
of the researcher, it is sometimes hard to filter out more relevant
from less relevant information, to understand the interplay of the
variables used in a study, and to interpret the results in a
theoretically meaningful way, beyond statements of the type ``variable
$x$ has a significant effect on variable $y$''.

While inferential statistics are indispensable in the comparative
analysis of corpora, it is important for researchers (and readers) to
also develop an intuitive understanding of the data (cf. also
\citealt{Griesetal2020}). Visualizations are a useful tool for this
purpose. In this chapter we discuss the use of specific tools for
visualizing bilingual data commonly encountered in contrastive
linguistics and translation studies (cf. also
\citealt{SoenningSchuetzler2023} on the visualization of corpus data
in general).

An obvious particularity of contrastive studies based on translation
corpora is the relational nature of the data, which often consists of
pairs of an original sentences and its translation or
interpretation. While the original sentence is an instance of
monolingual production, the translated sentence is determined both by
the co(n)text\footnote{`Cotext' is often used for the textual
  surrounding of a sentence and `context' for the extralinguistic
  circumstances of production. We will use `context' as a cover term
  for `cotext' and `context' narrowly defined in the following.}  and
the input from the source language. We thus distinguish between three
ways of analyzing frequency distributions. First, before carrying out
a comparative study we need to understand language-specific
distributions, using common tools of corpus linguistics; second, we
can compare distributions from different languages in a symmetric way
(`undirected comparison'); and third, we can study the choices made in
one language under the influence of another language,
e.g. translations of a given source language expression in a target
language (`directed comparison').

For illustration we use data analysed in previous studies by us
(\citealt{Atayanetal2019,GastAtayan2019,Gastetal2019,Atayan2021b,Atayan2021a})
on what we call adverbials of `immediate' or `proximal posteriority'. These
adverbials are introduced in Section \ref{sec:Background}, where we
also provide a brief description of the data.
%
In Section \ref{sec:LanguageSpecific}, we introduce techniques
commonly used for categorical (count) data in monolingual studies,
presenting some extensions for higher-dimensional data.
%
These techniques can also be used for undirected comparative data, as
is illustrated in Section \ref{sec:ComparativeUndirected}.
%
In Section \ref{sec:ComparativeDirected}, we demonstrate ways of
visualizing translation data. Sankey diagrams are useful for this
purpose, but have certain limitations. We therefore propose a new
type of visualization, Context-Conditional Translation Graphs, which
visualize translation choices made under specific contextual
conditions.
%
Section \ref{sec:ComparativeDirected} moreover contains some comments
on Bayesian networks, which can be used to factorize joint probability
distributions, and to visualize these factorizations, enabling
researchers to test models of bilingual language production. Section
\ref{sec:Conclusions} contains the conclusions.

%%% Local Variables:
%%% mode: latex
%%% TeX-master: "main"
%%% End:
